\documentclass{article}
\usepackage[utf8]{inputenc}
\usepackage{geometry}
\geometry{a4paper, margin=1in}

\title{Maturity as Orientation in Learning Networks}
\author{Albert Jan van Hoek}
\date{\today}

\begin{document}

\maketitle

Don’t be childish.

We say it easily.

But what do we actually mean?

As a child, you know you’re supposed to ``grow up.''
But when does that happen?

When can you say: last week I was still a child — this week I am an adult?

For a long time, I found that vague.

\section*{The TLC Perspective on Maturity}

The \textbf{Theory of Long-term Collaboration (TLC)} offers a surprisingly clear answer.

In a network of learning systems, there are only two basic orientations:
\begin{itemize}
    \item You act in ways that protect and strengthen the collaboration with others.
    \item Or you act in ways that protect yourself at the expense of that collaboration.
\end{itemize}
That’s the divide.

\section*{Mature vs. Childlike Behavior}

Mature behavior protects the long-term learning capacity of the network.
Childlike behavior protects the short-term ego of the individual.

Mature behavior:
\begin{itemize}
    \item Listens.
    \item Updates.
    \item Explains.
    \item Takes responsibility.
    \item Asks for help when needed.
\end{itemize}

Childlike behavior:
\begin{itemize}
    \item Hides.
    \item Deflects.
    \item Denies.
    \item Avoids ownership.
    \item Pretends to know.
\end{itemize}

And we all recognize the difference — not just in others, but in ourselves.
We know how it feels to take responsibility.
And we know how it feels to dodge it.

\section*{Maturity and Collaboration}

In that sense, maturity is not about age.
It is about orientation.
It is about having the nose in the same direction.

Long-term collaboration only happens when we all act maturely.

And therefore we need a network of adults.

We can’t afford to be childish.

\end{document}
